\cotefile{157-4}
\scanpage{82}
\mspage{75}
% archive sheet number: 81
et de plus \inserted{\unsure{avec}} $C,\mathrm{TC}$ \deleted{\unsure{de type fini}} tels qu'on ait
\[
C=\widetilde{C}_0,\qquad \mathrm{TC}=\widetilde{\mathrm{TC}}_0,
\]
avec $C_0,\mathrm{TC}_0$ petits. On pose maintenant
% slanted
\margin{On dit des $C,\mathrm{TC}$ (\unsure{en tant} que parties \unsure{closes} à droite) ``\unsure{accessibles|coaccessibles}''.}
\[
\begin{aligned}
F(I)&=\{f\in\mathrm{Fl}\,M(I)\mid f_s\in F\quad \forall s\in\mathrm{Ob}\,I\}\,,\\
\mathrm{TF}(I)&=\{f\in\mathrm{Fl}\,M(I)\mid f_s\in\mathrm{TF}\quad \forall s\in\mathrm{Ob}\,I\}\,,
\end{aligned}
\]
puis
\[
\begin{aligned}
C(I)&=\mathrm{TF}(I)^{\unsure{\alpha|\ast}},\\
\mathrm{TC}(I)&=F(I)^\alpha.
\end{aligned}
\]
L'espoir (que je devrai vérifier ultérieurement), c'est qu'on obtient ainsi une quillénisation uniforme à droite \inserted{de $W$,} sur $\mathrm{Cat}_{\unsure{\mathcal U|\mathscr U}}$ \unsure{tout entier}.

Le seul problème, c'est la validité des axiomes de factorisation de $f\in\mathrm{Fl}\,M(I)$ en $f=\unsure{pi|pj}$, avec $i\in\mathrm{TC}$, $\unsure{f|p}\in F$, et aussi avec $i\in C$, $f\in\mathrm{TF}$. J'ai bon espoir que c'est vrai --- cela se réduit essentiellement à voir pour $(C(I),\mathrm{TF}(I))$ \deleted{\unsure{soit}} et $(\mathrm{TC}(I),F(I))$, \deleted{\unsure{on doit le}} soit $(\unsure{\Phi|\varphi}(I),\unsure{\Psi|\Upsilon}(I))$, \deleted{\unsure{on doit prouver|on suppose}} en plus de $\Phi(I)=\Psi(I)^\alpha$, à $\Psi(I)=\Phi(I)_{\unsure{\alpha\delta|\alpha\beta}}$ (i.e. $\Psi(I)$ est clos à droite), et que $\Phi(I)$ soit accessible (tout comme $\Phi=\Phi(e)$).

Si ça marche, on aura l'\deleted{\unsure{admissibilité}}

\scanpage{83}
\mspage{76}
% archive sheet number: 82
des $u_!$ pour $u$ un foncteur entre cat. ordonnées finies, et sous des conditions largement plus générales. (\deleted{\unsure{Par}} P.ex. $u:I'\longrightarrow I$ avec $I$ ordonné, $I'$ ordonné et esp. sous-jacent noeth. et sobre.) Peut-être que cette fois, l'hypothèse que $W$ est stable par $\varinjlim$ filtrantes \unsure{va} quand même servir! Ainsi, le cas des catégories discrètes, et des catégories filtrantes à gauche comme $\mathbb N^{\unsure{\circ|o}}$, ne va plus poser de problèmes, à cause de cela. Il y a donc de l'espoir! Tout \deleted{catégorie} \inserted{ens.} ordonnée qui est $\varinjlim$ \unsure{filtrante|filtrantes} de sous-\deleted{\unsure{$I$|cat.}} \inserted{ens. ordonnés} \uline{fermés} noethériens et sobres, \unsure{est|\deleted{est}} \inserted{donc $W$-}admissible à droite. Mais qu'en est-il des catégories du type $I=B_G$ p.ex. ???

\scanpage{84}
\mspage{77}
% archive sheet number: 83
7. Filtrations cardinales et parties accessibles d'une \unsure{gross.|grosse} catégorie.

Je me réfère à la théorie des foncteurs et catégories accessibles de SGA 4, I, \S\ 9 (p. 138--178). Je me donne une grosse catégorie (disons une $\unsure{\mathcal U|U}$-catégorie) $E$ avec une \uline{filtration cardinale} (croissante)
\[
(\mathrm{Filt}^{\pi}(E))_{\pi\geq\pi_0},
\]
\margin{$\pi_0$ un card. infini fixé}
par des sous-catégories str. pleines $\mathrm{Filt}^{\pi}(E)$.

On \unsure{a|.}

a) $\mathrm{Filt}^{\pi}(E)$ \deleted{\unsure{est}} essentiellement petites, i.e. \unsure{équivalentes|équivalente} à une petite catégorie (\unsure{ou.|ou} à une catégorie $\in\mathcal U$)

b) $E$ satisfait $L_{\pi_0}$ (i.e. $E$ est stable par $\varinjlim$ de type $I$, quand $I$ est un ens. ordonné grand devant $\pi_0$), et $\mathrm{Filt}^{\pi}(E)$ est stable par les limites inductives filtrantes indexées par $I$ grand devant $\pi_0$ et tel que $\operatorname{card}I\leq\pi$.

c) Pour tt $\pi\geq\pi_0$ et pour tt $X\in\mathrm{Ob}\,E$, on a
\[
X\simeq\varinjlim_{i\in I}X_i
\]
où les $X_i$ ds $\mathrm{Filt}^{\pi}(E)$, et où $I$ est un ens. ordonné grand devant $\pi$.
% slanted
\margin{NB b) et c) impliquent que $E=\bigcup_{\pi\geq\pi_0}\mathrm{Filt}^{\pi}(E)$, \deleted{\uline{NB} \unsure{On a donc}}}
De plus, si $X\in\mathrm{Ob}\,\mathrm{Filt}^{\pi'}(E)$, avec $\pi'\unsure{\geq|>}\pi$, on peut prendre $I$ tel que $\operatorname{card}I\leq\pi'^{\pi}$ (donc si $\pi'=2^{c'}$ \inserted{avec $c'\geq\pi$}, d'où $(\pi')^{\unsure{n|\pi}}=(2^{c'})^{\pi}=2^{c'\pi}=2^{c'}=\pi'$, on aura $\operatorname{card}I\leq\pi'$. (\uline{NB} Si on admet l'hyp. du continu généralisée, on aura tjrs \deleted{$\operatorname{card}I\leq$} $\pi'^{\pi}=\pi'$ donc $\operatorname{card}I\leq\pi'$)

\scanpage{85}
\mspage{78}
% archive sheet number: 84
% Three long diagonal cancellation strokes cover the written draft.
\deleted{
J'ai \unsure{envie} ici de préciser la forme des systèmes inductifs intervenant dans c, quitte à faire donc une hyp. supplémentaire. Je vais considérer, \inserted{$\pi\geq\pi_0$ fixé, \unsure{et|tt}} pour \deleted{\unsure{$X$ fixé}} tt $X$, l'ens. ordonné $I_\pi(X)$ défini ainsi
\[
\begin{aligned}
I_\pi(X)={}&\text{ens. des sous-objets }\inserted{X_i}\text{ de }X\text{ tels}\\
&\text{que }X_i\in\mathrm{Filt}^{\pi}(E),\text{ ordonné par}\\
&\text{inclusion.}
\end{aligned}
\]
Il faut voir que cet ensemble \inserted{ordonné} est grand devant $\pi$, et que
\[
\varinjlim_{i\in I_\pi(X)}X_i=X.
\]
Si $X\in\mathrm{Ob}\,\mathrm{Filt}^{\pi}(E)$, alors $X$ est un plus grand \deleted{objet} élément de $I_\pi(X)$, qui est donc grand devant \unsure{tt|tout}, et $\varinjlim X_i=X_X=X$.

Considérons \deleted{\unsure{donc un|donc|le cas}} le cas général $X\in\mathrm{Ob}\,\mathrm{Filt}^{\pi'}(E)$, avec $\pi'>\pi$. Soit \deleted{\unsure{un|Soit}} $J\subset I_\pi(X)$ avec $J$ de cardinal $\leq\pi$. Il faut donc voir qu'il existe $Y\subset X$, \deleted{\unsure{avec}} avec $Y\in\mathrm{Filt}^{\pi}_{\deleted{\unsure{\pi|0}}}(E)$, et $Y\supset X_i$ pour tout $i\in J$.
}

\scanpage{86}
\mspage{79}
% archive sheet number: 85
Considérons \unsure{les|\deleted{les}} représentations de c)
\[
X=\varinjlim_I X_i\quad
\left\{\begin{aligned}
X_i&\in\deleted{\unsure{\mathrm{Ob}}}\mathrm{Filt}^{\pi}\\
I&\text{ grand devant }\pi
\end{aligned}\right.
\]
Soit
\[
\widetilde I=\{I'\subset I\mid\operatorname{card}I'\leq\pi,\quad I'\text{ grand devant }\pi_0\};
\]
\unsure{doué|muni} de la relation d'ordre d'inclusion.

Pour $I'\in\widetilde I$, soit $X_{I'}=\varinjlim_{I'}X_i$, cette limite existe par $L_{\pi_0}$ (Hyp. b), elle est dans $\mathrm{Filt}^{\pi}$ par b), donc on a un système inductif
\[
(X_{I'})_{I'\in\widetilde I},\qquad X_{I'}\in\mathrm{Filt}^{\pi},
\]
avec
\[
\operatorname{card}\widetilde I\leq(\operatorname{card}I)^{\pi}
\]
et si $X\in\mathrm{Filt}^{\pi'}$ $(\pi'>\pi)$, \deleted{\unsure{on peut|on prend}} prenant
\[
\operatorname{card}I\leq\pi'^{\pi}
\]
on aura
\[
\operatorname{card}\widetilde I\leq(\pi'^{\pi})^{\pi}=\pi'^{\pi^{\unsure{2|?}}}=\pi'^{\pi},
\]
donc
\[
\operatorname{card}\widetilde I\leq\pi'^{\pi}.
\]
Je dis que $\widetilde I$ est lui aussi \uline{grand devant} $\pi$, \unsure{tout|tt} comme $I$. En effet,

\scanpage{87}
\mspage{80}
% archive sheet number: 86
soit
\[
(I_\alpha)_{\alpha\in A},\quad I_\alpha\in\widetilde I,\quad\operatorname{card}A\leq\pi
\]
\deleted{\unsure{et considérons}} prouvons qu'il existe $I'\in\widetilde I$ qui majore les $I_\alpha$, i.e. qui contient $I'_0=\bigcup_{\alpha\in A}I_\alpha$. On aura $\operatorname{card}I'_0\leq\pi^2=\pi$ (car $\operatorname{card}A\leq\pi$ et $\operatorname{card}I_\alpha\leq\pi$ pour tt $\alpha$), et il suffit de voir ceci:

\uline{Lemme 1}: Soit \inserted{$\pi_0,\pi$ deux cardinaux, avec $\pi\geq\pi_0$,} $I$ un ens. ordonné grand devant $\pi$, et $I'_0$ une partie de $I$ de cardinal $\leq\pi_0$. Alors il existe une partie $I'$ de $I$, \deleted{contenant $I'_0$, et} satisfaisant les conditions
\[
\begin{aligned}
\text{a)}\quad&I'\supset I'_0\\
\text{b)}\quad&\operatorname{card}I'\leq\pi^{\ast}\,,\\
\text{c)}\quad&I'\text{ grand devant }\pi_0.
\end{aligned}
\]
% slanted; linked to the star beside the cardinal bound.
\margin{Pourvu qu'on suppose $\pi>\pi_0$, et de plus $\pi^{\pi_0}=\pi$ (p.ex. $\pi=2^c$ avec $c\geq\pi_0$)}
On se propose de construire $I'$ par récurrence transfinie, en définissant les $I'_\alpha\subset I^{\unsure{\deleted{\xi}|\xi|\deleted{\ast}}}$ ainsi:
\[
\begin{aligned}
I'_0&=I'_0\\
I'_{\alpha\deleted{\unsure{a|0}}}&=\bigcup_{\beta<\alpha}I'_\beta\quad\text{si }\alpha\text{ est un ordinal limite}
\end{aligned}
\]
$I'_{\alpha+1}$ déduit de $I'_\alpha$ ainsi: pour tte partie $J$ de $I'_\alpha$, telle que $\operatorname{card}J\leq\pi_0$, qui \unsure{n'est pas majorée} dans $I'_\alpha$, on choisit un \unsure{major.|majorant} \inserted{$i_J$} dans $I$ (existe, car $I$ grand devant $\pi$ donc dev. $\pi_0$), et $I'_{\alpha+1}$ \unsure{s'obtient} en adjoignant ces $i_J$ à $I'_\alpha$.

\scanpage{88}
\mspage{81}
% archive sheet number: 87
Soit alors $\alpha_0$ le premier ordinal qui soit de cardinal $\deleted{\unsure{\geq|\leq}}>\pi_0$, et considérons \deleted{$\unsure{I|I'}_{\alpha_0}$}
\[
I'=I'_{\alpha_0}=\bigcup_{\substack{\alpha<\alpha_0\\\text{i.e. }\operatorname{card}\alpha\leq\pi_0}}I'_\alpha.
\]
Alors $I'$ contient $I'_0$, il est grand devant $\pi_0$, mais \inserted{quel est} son cardinal. Prouvons par récurrence transfinie que
\[
\operatorname{card}I'_\alpha\leq\pi\qquad
\begin{gathered}
\text{si }\alpha<\alpha_0\\
\text{i.e. }\operatorname{card}\alpha\leq\pi_0.
\end{gathered}
\]
C'est OK pour $\alpha=0$ par hyp. sur $I'_0$. Si c'est vrai pour $I'_\alpha$, est-ce vrai pour $I'_{\alpha+1}$? On \deleted{lui} adjoint à $I'_\alpha$ un ens \unsure{d'éléments|d'élémts} $\unsure{\deleted{e}i_J|i_J}$ qui est de cardinal $\leq\unsure{\operatorname{Card}|\overline{\mathrm C}\mathrm{ard}}\unsure{\mathscr P|\mathcal P|P}_{\leq\pi_0}(I'_\alpha)$ donc de cardinal $\leq\operatorname{card}(I'_\alpha)^{\pi_0}\leq\pi^{\pi_0}$, et si on suppose
\[
\pi^{\pi_0}=\pi,
\]
on \unsure{gagne|gagné}, \deleted{\unsure{Cela sera le cas}} Supposons enfin que $\alpha$ soit un ordinal limite (de cardinal $\leq\pi_0$). Donc on a
\[
I'_\alpha=\bigcup_{\beta<\alpha}I'_\beta
\]
d'où \deleted{$\operatorname{card}$}

\scanpage{89}
\mspage{82}
% archive sheet number: 88
\[
\operatorname{card}I'_\alpha\leq
\underbrace{(\operatorname{card}\alpha)\pi}_{\leq\pi.\pi=\pi^2=\pi}
\qquad\text{puisque les }\operatorname{card}I'_\beta\leq\pi,
\]
on gagne encore.

On trouve enfin, par les $=$ arguments
\[
\operatorname{card}I'_{\alpha_0}\leq(\operatorname{card}\alpha_0)
\underbrace{\pi\leq\pi^2}_{\substack{\text{si }\operatorname{card}\alpha_0\leq\pi\text{ i.e.}\\\text{si }\pi_0<\pi}}=\pi,
\]
donc on gagne encore.

Ainsi, on a trouvé un deuxième système inductif
\[
(X_{\widetilde i})_{\widetilde i\in\widetilde I}
\]
de $=$ type que le système initial, ayant comme limite inductive $X$ (qu'il ait $=$ limite inductive est immédiat). Mais qu'a-t-on gagné? Ceci: \unsure{Dans|\deleted{Dans}} \deleted{$\unsure{I|X}$} $\widetilde I$ \unsure{tte|toute} \deleted{\unsure{Toute partie ordonnée|Toute partie}} partie \inserted{de $\widetilde I$ et \unsure{grande|grand} devant $\pi_0$} de cardinal $\leq\pi$ est non seulement \uline{majorée} dans $\widetilde I$, mais même y admet une borne supérieure.
\margin{$\Vert$}
\deleted{En effet, si on a des $I'_\alpha\in\widetilde I$, i.e. des $I'_\alpha\subset I$ qui}
\inserted{Cela résulte des}

\scanpage{90}
\mspage{83}
% archive sheet number: 89
\uline{Lemme 2}: Soit $I$ grand devant $\pi>\pi_0$, avec $\pi^{\pi_0}=\pi$ (cf lemme 1). Soient $(\unsure{I'|E'})_{\alpha\in A}$ une \uline{famille} de parties de $I$ qui sont dans $\widetilde I$ i.e. a) de cardinal $\leq\pi$ et b) grandes devant $\pi_0$; et supposons que \inserted{$\operatorname{card}A\leq\pi$ et que} $A$, ordonné par inclusion, soit grand devant $\pi_0$. Alors $I'=\bigcup_{\alpha\in A}I'_\alpha$ est encore dans $\widetilde I$, i.e. $I'$ est de card $\leq\pi$ (car de card $\leq\pi^2=\pi$) et de plus, grand devant $\pi_0$. (Vérification immédiate)

Ainsi, on a prouvé ceci

\uline{Proposition 1} Soit $E$ grosse catégorie, avec filtration cardinale $(\mathrm{Filt}^{\pi})_{\pi\geq\pi_0}$. Soit $\pi$ un cardinal avec
\margin{$\big|$}
\[
\pi>\pi_0,\qquad\pi^{\pi_0}=\pi\qquad(\text{p.ex. }\pi=2^c,\text{ avec }c\geq\pi_0).
\]
Alors pour tt $X$ dans $E$, on peut trouver un système inductif $(\unsure{K|X}_i)_{i\in I}$, \deleted{\unsure{de $=$ limite $X$, $I$ grand dev.}} \inserted{de limite $X$,}
\margin{$\big|$}
\[
\text{a)}\quad\left\{\begin{aligned}
X_i&\in\mathrm{Filt}^{\pi}\ \deleted{\unsure{E|I}}\quad\text{pour tt }i\in I\\
I&\text{ grand devant }\pi,\\
&\deleted{\unsure{\varinjlim X_i\simeq X}}
\end{aligned}\right.
\]
satisfaisant aux conditions supplémentaires suivantes

b) Toute partie $\unsure{J|\mathcal J}$ de $I$ de cardinal $\leq\pi_0$, et grande devant $\pi_0$, admet une borne supérieure $i_J$ dans $I$.

\unsure{c)|b)} On \deleted{de cette borne supérieure} a, dans les conditions

\scanpage{91}
\mspage{84}
% archive sheet number: 90
\[
X_{i_J}\xleftarrow{\unsure{\sim|\simeq}}\varinjlim_{i\in J}X_i
\]
De plus, si $X\in\mathrm{Filt}^{\pi'}$ avec $\pi'>\pi$, on peut prendre $I$ de cardinal $\leq\pi'^\pi$ (donc de cardinal $\leq\pi'$ si $\pi'^\pi=\pi'$, p.ex. si $\pi'=2^{c'}$, avec $c'\geq\pi$).

\uline{Définition} On appelle un \deleted{\unsure{représentation}} \inserted{système} inductif $(X_i)_{i\in I}$ satisfaisant les conditions a), b), c) ci-dessus, un \uline{système inductif bien adapté} (à la filtration cardinale donnée, pour les \unsure{cardinaux|cardinalités} \deleted{\unsure{$\pi_0$}} $\pi_0,\pi$ \deleted{\unsure{fixés|fixées}} \unsure{donnés|données}), et la représentation correspondante
\[
X\xleftarrow{\sim}\varinjlim X_i
\]
une \uline{représentation bien adaptée} de $X$ en termes d'objets de $\mathrm{Filt}^{\pi}$.

\uline{NB} Quand $E$ est stable par petites $\unsure{\varinjlim|\underline{\lim}}$ filtrantes, la situation se simplifie considérablement. On peut alors prendre
\[
\widetilde I=\{I'\subset I\mid\operatorname{Card}I'\leq\pi,\ I'\text{ filtrant}\}
\]
(sans \unsure{s'astreindre à|s'encombrer de} la condition $I'$ grand devant $\pi_0$) et définir encore le système inductif $(X_{I'})_{I'\in\widetilde I}$ par
\[
X_{I'}=\unsure{\lim|\varinjlim}_{i\in I'}X_i.
\]
On a \unsure{encore|alors}

\scanpage{92}
\mspage{85}
% archive sheet number: 91
\[
\operatorname{card}\widetilde I\leq(\unsure{\operatorname{Card}|\operatorname{card}}I)^{\pi}\leq\pi'^\pi\qquad\text{si }\operatorname{card}I\leq\pi'^\pi
\]
(notamment, si $X\in\mathrm{Filt}^{\pi'}$ et qu'on prenne $I$ comme indiqué dans la condition c des filtr. cardinales). De plus, cette fois, le fait que $\widetilde I$ soit grand devant $\pi$ est trivial, on a même mieux, \inserted{\unsure{: pourvu que} $I$ soit filtrant:} car \uline{Toute} partie \inserted{$(I'_\alpha)_{\alpha\in A}$} de $\widetilde I$ \inserted{filtrante et} de cardinal $\leq\pi$ admet une borne supérieure \deleted{\unsure{$I$}} $\tilde i_A$ (qui est bien filtrante).

savoir la \inserted{réunion} \deleted{\unsure{partie engendrée et|partie engendrée|partie} filtrante de $I$ engendrée par \inserted{\unsure{on prend} $I'_0$} $I'_0=\bigcup_{\alpha\in A}I'_\alpha$, \unsure{qui est une} partie de cardinal $\leq\pi^2=\pi$, et qui \unsure{la|est une} partie filtrante}

De plus, le système inductif $(\unsure{X|\alpha}_{\tilde i})_{\tilde i\in\widetilde I}$ a la propriété plus forte que
\[
X_{\tilde i_A}\xleftarrow{\sim}\varprojlim_{\tilde i\in A}X_{\tilde i}.
\]
Ainsi on a la

\uline{Proposition 1 bis} Supposons \deleted{$\mathrm{Filt}^{\pi}$ stable} \inserted{$E$ stable par petites} $\unsure{\varinjlim|\underline{\lim}}$ filtrantes, et $\mathrm{Filt}^{\pi}$ stable par celles-ci. \deleted{Alors} si $\pi\geq\pi_0$ est donné. Alors pour tt $X$ dans $E$, on peut trouver un syst. ind. $(X_i)_{i\in I}$ de $\varinjlim_I X$, avec

a) $X_i\in\mathrm{Filt}^{\pi}$ pour tt $i\in I$, $I$ ens. ordonné grd dev. à $\pi$

b) Toute partie $J$ de $I$ filtrante de card. $\leq\pi$ admet une borne sup. $i_J$ dans $I$,

c) \unsure{dans c)|dans ce}, sous les conditions

\scanpage{93}
\mspage{86}
% archive sheet number: 92
\[
X_{i_J}=\varinjlim_{i\in J}X_i.
\]
De tels syst. inductifs seront dits ``\uline{parfaitement adaptés}'' (à la filtration cardinale donnée, pour le cardinal $\pi$ donné --- $\pi_0$ n'y figure plus, si ce n'est comme début de la filtr. cardinale en \unsure{place|plus}.).

Comme pratiquement, on aura à considérer des $E$ stables par \inserted{petites} $\varinjlim$ filtrantes, et des filtrations cardinales stables \deleted{\unsure{par}} \unsure{et coniques|et canoniques} (renforcement de la condition b) des filtr. cardinales), c'est le cas, \deleted{\unsure{les}} et la notion des représentations \unsure{parfaitement|bien} adaptées de $X$ comme $\varinjlim_I\unsure{X_i|\alpha_i}$ ($X_i\in\mathrm{Filt}^{\pi}$), que je mettrai en avant. Le cas ``général'', où on suppose simplement $E$ satisfait $L_{\pi_0}$, est plus alambiqué et se traite moins joliment.
% slanted
\margin{On parlera de filtration cardinale \uline{stricte}, si les \uline{$\mathrm{Filt}^{\pi}$} sont stables par $\varinjlim$ filtrantes de cardinal $\leq\pi$. (et $E$ par les limites \unsure{filtrantes} petites)}

\uline{Définition 2} Soit $E$ muni d'une filtration cardinale stricte $(\underline{\mathrm{Filt}^{\pi}})_{\pi\geq\pi_0}$. \deleted{Une} \inserted{Soit $M$ une} partie \deleted{\unsure{$M$}} de $\operatorname{Ob}E$, $\underline M$ la sous-cat. pleine de $E$ qu'elle définit, \inserted{et soit $\pi$ un cardinal \deleted{\unsure{$\pi$}} $\geq\pi_0$.} On dit que \unsure{$M$} est \uline{$\pi$-accessible}, ou que $\underline M$ est $\pi$-accessible,

\scanpage{94}
\mspage{87}
% archive sheet number: 93
si les conditions suivantes sont satisfaites.

Pour tt \deleted{\unsure{$X\in\operatorname{Ob}E$, et}} système inductif $(X_i)_{i\in I}$ parfaitement adapté à la filtration cardinale pour $\pi$ (donc $I$ ens. ordonné \deleted{grand devant $\pi$;} \inserted{filtrant} tel que tt partie filtrante $\unsure{J|I}$ de cardinal $\leq\pi$ ait une borne supérieure $i_J$, et $X_{i_J}\xleftarrow{\simeq}\varinjlim_{j\in J}X_j$), \deleted{de limite inductive $X$ \uline{ds $M$},} \deleted{\unsure{alors pour tt $i\in I$ on ait}.} \inserted{Soit $I_M$ l'ens des} \deleted{\unsure{l'ens}} \deleted{$I_M$ des} $i\in I$ tels que $X_i\in M$. \deleted{est cofinal dans $I$} Ceci posé, on veut ceci:

a) pour que $X$ ($=\varinjlim X_i$) soit dans $M$, il \unsure{faut (et) suffit|faut et suffit}, \deleted{\unsure{par définition}} \deleted{\unsure{que}} \inserted{que} $I_M$ soit cofinal dans $I$ (de sorte que $X$ apparaisse comme $\varinjlim$ filtrante, d'un ens. d'indices $I_M$ grand devant $\pi$, d'éléments de $M\cap\mathrm{Filt}^{\pi}\overset{\mathrm{def}}{=}M^{\pi}$).

\deleted{b) $M$ est stable par \deleted{petites} limites \unsure{inductives} filtrantes \inserted{(\unsure{en $\underline M$}),} grandes devant $\pi$.
% slanted
\margin{\unsure{Découle|Donc} \uline{donc} la condition b), qui est indépendante d'un syst. inductif particulier.}

On dit que $M$ est accessible s'il existe un cardinal $\pi\geq\pi_0$ tel que $M$ soit $\pi$-accessible.}

\scanpage{95}
\mspage{\unsure{88|89}}
% archive sheet number: 94 (upper left)
\deleted{$\unsure{J|I}$ de cardinal $\leq\pi$ et grand devant $\pi$, soit $I$}

J'ai envie de reprendre la définition de l'a\deleted{\unsure{cc}} $\pi$-accessibilité de $M$, pour une filtration cardinale donnée à coups de conditions équivalentes. Pour tt système inductif \deleted{\textcircled{$\unsure{a|\alpha}$}} $(X_i)_{i\in I}$ dans $M$, je garde la notation
\[
I_M=\{i\in I\mid X_i\in M\}
\]
Je pose, pour tt cardinal $\pi\geq\pi_0$
\[
M^{\pi}=M\cap\mathrm{Filt}^{\pi}
\]
\deleted{\textcircled{$\unsure{a_\pi|b}$} $M$}

\textcircled{$a_\pi$} $\unsure{M|\underline M}$ est l'ens des \deleted{\unsure{limites}} objets de la forme
\[
\varinjlim_{i\in I}X_i,
\]
où les $X_i$ sont dans $M^{\pi}$, et $I$ est grand devant $\pi$.

\textcircled{\textcircled{$b_\pi$}} Soit $(X_i)_{i\in I}$ un syst. inductif dans $M$, \deleted{$I$ grand devant $\pi$, et} \inserted{parfaitement adapté à la filtration cardinale, pour $\pi$, i.e. $I$ filtr., et toute partie filtrante $J$ de cardinal $\leq\pi$ admet une borne supérieure $i_J$, et $X_{i_J}\xleftarrow{\simeq}\varinjlim_J X_i$, et enfin $X_i\in\mathrm{Filt}^{\pi}$ pour tt $i\in I$.} Soit
\[
X=\varinjlim X_i.
\]
Pour que l'on ait $X\in M$, il f. et s. que $\unsure{I_M|I^{\pi}_M}$ soit cofinal dans $I$.

\textcircled{$c_\pi$} $\underline M^{\pi}=\underline M\cap\underline{\mathrm{Filt}}^{\pi}$ est une souscatégorie de $\underline M$ génératrice par épimorphismes stricts,

\textcircled{$d_\pi$} \deleted{[et par \unsure{limites}]} $M$ est stable par $\varinjlim$ grandes devant $\pi$, enfin pour tt $X\in M$, $\underline M^{\pi}/X$ est une catégorie filtrante.$^\ast$
% slanted
\margin{$\ast$ cf plus bas pour une \unsure{vérif.}}

\scanpage{96}
\mspage{89}
% archive sheet number: 95
\deleted{\textcircled{$\unsure{e'|e}$} $\underline M$ satisfait $L_\pi$, \deleted{\unsure{et}} le foncteur d'inclusion $\underline M\longrightarrow E$ est $\pi$-accessible $^\ast$[et admet \deleted{\textcircled{$\unsure{d|e|c}$}} une sous-catégorie strictement génératrice contenue dans $\underline M^{\pi}$].
% slanted
\margin{Cette condition \unsure{semble plus} \unsure{impliquée|compliquée}, \unsure{simplement|sûrement}.}}

Je présume que ces conditions sont équivalentes.

$b_\pi\Rightarrow a_\pi$. On suppose b) vérifié, on veut prouver a. D'abord que \deleted{tout} $X\in M$ est de la forme dite $\varinjlim_I X_i$, avec les $\unsure{X_i|\alpha_i}\in M^{\pi}$ et $I$ grand devant $\pi$. Par prop. 1', on sait qu'on a une présentation parfaitement adaptée pour $\pi$, $X=\unsure{\underline{\lim}|\varinjlim}_{i\in I}X_i$. Par b), \inserted{``il faut''} la partie $I_M$ de $I$ est cofinale à $I$, donc comme $I$ est grand devant $\pi$, $I_M$ aussi, et on aura (posant $J=I_M$)
\[
X=\varinjlim_{i\in I_M}X_i,\qquad\begin{aligned}
X_i&\in M^{\pi},\\
I_M&\text{ grand devant }\pi,
\end{aligned}
\]
OK. Inversement, prouvons que si $X$ se met sous la forme dite, \inserted{dans a),} il est dans $M$. Cela va résulter du

\uline{Lemme 1}: Si \deleted{\unsure{les}} la condition \deleted{\unsure{a)}} b) \deleted{\unsure{(on admet $a_\pi$, qui est impliqué par b))}} est satisfaite, $M$ est stable par limites inductives grandes devant \deleted{\unsure{$\pi_0$|$<$}} $\pi$.
% slanted
\margin{\deleted{\unsure{On peut} \unsure{prendre} \unsure{filtrant}}}

\deleted{\unsure{Pour le moment, on a}}

Il faut prouver que si $\unsure{\deleted{X}X_i|X_i|X}$

\scanpage{97}
\mspage{\unsure{90}}
% archive sheet number: 96
\deleted{\unsure{des $X_i$ quelconques}}
\[
X=\varinjlim_{i\in I}X_i,
\]
$I$ grand devant $\pi$, les $X_i$ dans $M$, alors $X\in M$. La condition b) \inserted{``il suffit''} l'assure$^\ast$ quand on suppose de plus les $X_i$ dans $\mathrm{Filt}^{\pi}$, (donc dans $M^{\pi}$), de sorte que l'on aura $I_M=I$, donc $I_M$ cofinal, donc $X\in M$.
% slanted
\margin{pas bien: faut \unsure{encore} supposer plus, sur le système inductif. On ne peut \unsure{tjrs} disposer \unsure{des} \unsure{lim.} \unsure{en ce cas}, $X_i\in M^{\pi}$ $\forall i\in I$, \unsure{le|et} \unsure{faire}, \unsure{le} \unsure{raisonnement} \unsure{simplement}.}
% slanted
\margin{\deleted{\uline{NB} \unsure{ce qu'il faut} \unsure{faire} \illegible{2} pour $b\Rightarrow a$.}}

Si on ne suppose pas les $X_i$ dans $\mathrm{Filt}^{\pi}$, il faut travailler un peu plus, en \deleted{prouvant ceci:}
\margin{Passer \unsure{ici} de suite à la partie de la démonstration ci-dessous page 96 (qui prouve $b\Rightarrow a$). \unsure{Le reste de la démonstr.} \unsure{n'est pas prouvé}.}
% slanted
\margin{\uuline{Ultra-faux!!}}
\deleted{\uline{!! Lemme 2} La condition b) \unsure{implique}

La condition b) \unsure{s'implique|s'entend}, plus fort, si on n'exige plus que les $X_i$ soient dans $\mathrm{Filt}^{\pi}$.

Soit en effet $x=(X_i)_{i\in I}$ avec $I$ \inserted{filtrant et} tel que tt partie filtr. $J$ de cardinal $\leq\pi$ ait une borne supérieure $i_J$, et qu'on ait $X_{i_J}\xleftarrow{\unsure{\sim|\simeq}}\varinjlim_{i\in J}X_i$. Pour chaque $i\in I$,

il existe par prop. 1', une présentation parfaite (pour $\pi$) de $X_i$,
\[
X_i=\varinjlim_{\alpha\in I_i}X_{i,\alpha}
\qquad\left\{\begin{aligned}
&\text{les }X_{i,\alpha}\in\mathrm{Filt}^{\pi}\\
&I_i\text{ grand devant }\pi\\
&\text{et de plus, tt }\unsure{A}\text{ partie de }I_i\\
&\text{cardinal }\leq\pi\text{ a une borne supérieure}\\
&\alpha_A,\text{ etc}
\end{aligned}\right.
\]}

\uline{Lemme 2} Si tout élément de $\mathrm{Filt}^{\pi}$ est $\pi$-admissible, alors $b_\pi)$ (avec $a_\pi$) implique que $M$ est stable par $\varinjlim$ grandes devant $\pi$.
